\chapter{GIỚI THIỆU CHUNG}

\section{Tổng quan đồ án}

Đồ án \textit{Internet Routing} thuộc học phần Mạng máy tính nâng cao,
tập trung mô phỏng hoạt động định tuyến giữa các hệ tự trị (AS)
trên môi trường GNS3. Mô hình gồm ba AS A, B và C,
sử dụng RIPv2, OSPF và BGP để trao đổi thông tin định tuyến
và thiết lập kết nối giữa các mạng LAN.

\section{Mục tiêu và phạm vi thực hiện}

Đồ án giúp củng cố kiến thức về định tuyến nội miền, liên miền
và cách áp dụng chính sách lựa chọn đường đi. Các nội dung chính gồm:
\begin{itemize}
    \item Quy hoạch địa chỉ IP, cấu hình RIPv2, OSPF và các phiên iBGP, eBGP.
    \item Thiết lập chính sách BGP để ưu tiên đường đi và kiểm soát quảng bá tuyến.
    \item Thử lỗi liên kết, đo thời gian khôi phục kết nối
    và đánh giá ảnh hưởng của việc điều chỉnh timer BGP.
    \item Sử dụng Wireshark để phân tích các bản tin BGP và RIPv2.
\end{itemize}

\section{Phương pháp thực hiện}

Nhóm triển khai cấu hình theo từng giao thức, kiểm tra bảng định tuyến
và trạng thái láng giềng trước khi thử nghiệm các tình huống lỗi.
Kết quả được đối chiếu qua Ping, Traceroute và các gói tin bắt bằng Wireshark,
làm cơ sở đánh giá khả năng kết nối, lựa chọn tuyến và hội tụ của mạng.
